\section{LLM-Assisted Assembly Transformation}\label{sec:llm}
\subsection{What the Model Is Asked to Change}
In transformation mode, the LLM receives an assembly representation and
proposes a new instruction sequence intended to preserve the program's
behavior. This differs from Section~\ref{sec:code_data}, where its answer
contains only data locations. Zone supports four transformation requests:
equal code substitution (EQ) replaces instruction sequences; NOP insertion
adds operations intended to leave the relevant state unchanged; bogus
control flow (BCF) adds control-flow structure without changing the intended
outcome; and instruction block permutation (IBP) changes the placement of
blocks while retaining their execution relationships. These are the
families introduced in Section~\ref{sec:background}; the LLM proposes their
application rather than introducing a fifth family. PII and assembly
wrapping remain native engine operations. This is an implementation scope
choice: the LLM interface defines requests and validators for the four
listed families. PII and wrapping also impose constraints on the binary
layout and encoding that this interface does not implement. Their absence
from LLM mode is not a claim that a model could never propose them.

Before sending a request, Zone converts the binary into a prepared
assembly listing. Embedded data that has been identified and extracted is
kept separately, so the model proposes changes to the instruction
representation while the pipeline retains the exact data contents.
Identification may itself use an LLM, but copying retained data and checking
its contents are deterministic operations. This separation makes the role
being evaluated explicit. An end-to-end model that also rewrites data would
be a different design and is not compared in these experiments.
Relevant locations are named with symbolic labels. A label is a name for
a destination, like a section name in a document: it stays meaningful even
when edits change the destination's numerical position.

The model returns a complete proposed assembly listing. Zone then converts
that text back into binary instructions, restores the preserved data, and
performs its structural checks. There are therefore three successive
representations: the original binary, the prepared input listing, and the
final binary. The LLM similarity score compares the prepared and returned
listings; changes already made during preparation are outside that comparison.

The experiments request \texttt{deepseek-flash}. The locally recorded
version name comes from the client's alias table, rather than independent
provider confirmation, so the requested identifier is the reproducible
model designation available in these records. Both LLM roles use a
requested temperature of 0.2; the client supplies neither a sampling seed
nor an explicit top-p setting. Repeated requests are therefore not treated
as deterministic reproductions.

\subsection{Prompting Conditions and Repeated Attempts}
A \emph{zero-shot} request describes the transformation and its constraints
without a worked example. A \emph{two-shot} request adds two examples of
input and transformed output. The number of shots thus counts examples
inside a prompt; it does not count generated variants or execution tests.
Each input, transformation family, and shot setting defines an experimental
\emph{cell}.

If structural validation rejects an answer, the pipeline can send the
diagnostic back to the model and ask for a corrected proposal. This is
\emph{validation feedback}: an explanation of what the checker rejected,
such as assembly that cannot be built or a missing label. It is not a
functional success signal or a scanner verdict. A retry can therefore use
information unavailable in the initial request.
The benign benchmark disables
this mechanism and performs three separate, single-generation invocations
per cell, retaining every failure. The full-payload campaign instead
targets one output per cell, with some retries and restarted failed runs
present in the saved logs. Its results are reported as cell outcomes, not
as first-generation success rates. Section~\ref{sec:methodology} gives the
exact accounting for both protocols.

\subsection{From a Proposed Listing to an Assessed Output}
Validation checks whether the response can be assembled and decoded,
whether tracked control-flow targets remain consistent, whether preserved
data survives, and whether the requested transformation's structural
requirements are met. A response can fail at any of these stages, or be
truncated before it contains a complete listing. Passing them means that
Zone accepts a constructed output, which then proceeds to the execution test.

Functional assessment then uses an external \emph{test harness}, meaning
the program that launches a test and records its observable result. For
the controlled programs, it compares a complete output buffer with the expected
bytes fixed before generation. For the original payloads, it checks the
class-specific observable described in Section~\ref{sec:methodology}.
The functional test is applied after generation in these experiments. Saved records distinguish model requests, accepted outputs,
execution results, time, and estimated cost, allowing rejection and execution
failure to be analyzed independently.
